\bookchapter{פרק 56: ליטוש והוצאה לאור – בדיקות, ניפוי באגים ועלייה ל-Google Play}

\chapterintro{האפליקציה נבנתה – ומה עכשיו? הפרק הזה עונה על שלושה עניינים בדרך מ“זה רץ” אל “זה מוכן למסירה”: איך לנפות באגים ביעילות, איך לכתוב בדיקות שמונעות נסיגות, ואיך להעלות את האפליקציה ל-Google Play כדי שכל העולם יוכל להוריד אותה. עלייה לחנות היא לא רק טקס – זהו קו הפרדה בין יצירה למוצר.}

\booksection{א' שלושת כלי הניפוי: Logcat, נקודות עצירה ובודק ה-UI}

\begin{center}
\begin{tabularx}{\linewidth}{|X|X|X|}
\hline
\textbf{כלי} & \textbf{למה הוא משמש} & \textbf{איך משתמשים} \tabularnewline
\hline
Logcat & צפייה בלוגים של המערכת או שלכם & בקוד: Log.d("TAG","ערך=\$x"), ובפאנל מחפשים את ה-TAG לסינון \tabularnewline
ניפוי עם נקודות עצירה & לעצור את התוכנית בשורה מסוימת ולבחון ערכי משתנים & לוחצים ליד מספר השורה – מופיעה נקודה אדומה; מריצים במצב Debug ומתקדמים שורה אחר שורה \tabularnewline
Layout Inspector & צפייה במבנה המסך בזמן ריצה & בתפריט Tools Layout Inspector, צפייה בזמן אמת בעץ הרכיבים \tabularnewline
\hline
\end{tabularx}
\end{center}

\begin{center}\small\color{gray}שלושת כלי הניפוי\end{center}

\begin{tcbinfo}{טיפ: איך קוראים הודעת שגיאה?}
אחרי קריסה תופיע ב-Logcat מחסנית אדומה (Stacktrace); חפשו מלמטה למעלה את השורה הראשונה שמתחילה בשם החבילה שלכם (למשל com.example.yourapp) – שורה זו היא בדרך כלל זירת הפשע. קודם בדקו את “סוג ה-Exception” (מצביע ריק? המרת טיפוס?), ואז את “at ... הקובץ שלכם:מספר השורה”. הטריק הזה פותר 80\% מהקריסות.

\end{tcbinfo}

\booksection{ב' בדיקות: ביטוח לקוד}

בדיקה ידנית אינה אמינה – מתקנים תכונה אחת ואחרת נשברת בשקט. בדיקות הן “כדור התחרפות” וגם “שריון יורה”. בדיקות האנדרואיד מתחלקות לשלוש שכבות, ומתחילים צריכים קודם כול לשלוט בשכבה הראשונה:

\begin{center}
\begin{tabularx}{\linewidth}{|X|X|X|X|}
\hline
\textbf{שכבה} & \textbf{מה בודקים} & \textbf{היכן ממוקם} & \textbf{מהירות} \tabularnewline
\hline
בדיקות יחידה & לוגיקה טהורה (חישובים, ולידציה) & app/src/test/ & שניות \tabularnewline
בדיקות מכשיר (Instrumented) & התנהגות המסך על מכשיר אמיתי & app/src/androidTest/ & דקות \tabularnewline
בדיקה ידנית & תחושה וחוויה של אדם אמיתי & —— & תמיד נחוצה \tabularnewline
\hline
\end{tabularx}
\end{center}

\begin{center}\small\color{gray}שלוש שכבות הבדיקות\end{center}

\begin{lstlisting}[language=Kotlin]
// לחלץ את הלוגיקה של "המרת טקסט הסכום למספר" לפונקציה טהורה
fun parseAmount(input: String): Double? {
    val v = input.trim().toDoubleOrNull()
    return if (v != null && v > 0) v else null
}

// בדקו אותה (JUnit4)
@Test
fun emptyInputReturnsNull() {
    assertEquals(null, parseAmount(""))
}
@Test
fun negativeReturnsNull() {
    assertEquals(null, parseAmount("-5"))
}
@Test
fun normalValueParses() {
    assertEquals(12.5, parseAmount("12.5"))
}
\end{lstlisting}

את בדיקות הסכום והחישובים הסטטיסטיים בספר החשבונות צריך לכסות בבדיקות יחידה. הרצת בדיקות: קליק ימני על תיקיית test Run. ירוק – הכול עבר; עכשיו אפשר לשנות קוד בראש שקט.

\booksection{ג' חתימה ואריזה: לתת לאפליקציה “חותמת”}

אנדרואיד דורשת שכל חבילת התקנה תהיה חתומה – כמו חותמת רשמית שמאשרת “החבילה הזאת יצאה ממני ולא שונתה”. בזמן פיתוח Android Studio חותמת אוטומטית בחתימת debug; להפצה דרושה חתימת release משלכם. התהליך:

\begin{enumerate}
  \item בתפריט Build Generate Signed Bundle / APK;
  \item יוצרים קובץ keystore (מחסן מפתחות) – מגדירים סיסמה ושומרים עליה היטב;
  \item בוחרים release בפורמט AAB – והבנייה הושלמה.
\end{enumerate}

\textbf{אזהרה: keystore שאבד = אפליקציה שאבדה}\par

עדכון האפליקציה אחרי העלייה לחנות חייב להתבצע עם אותה חתימה, אחרת החנות תדחה את העדכון. אם קובץ ה-keystore והסיסמה אבדו, תיאלצו להעלות אפליקציה חדשה מאפס – וכל המשתמשים הוותיקים יישארו מחוץ למשחק. הקפידו לגבות את ה-keystore והסיסמה במקום בטוח (ענן מוצפן + מחברת רשומה) – זהו אחד השיעורים היקרים ביותר של מפתחי אנדרואיד.

\booksection{ד' התהליך המלא של עלייה ל-Google Play}

רישום חשבון מפתח

הכנת חומרים לחנות

העלאת ה-AAB

מילוי טופס אבטחת המידע

בדיקה (שעות עד ימים)

פרסום

\begin{center}\small\color{gray}איור 6: שישה שלבים לעלייה ל-Google Play\end{center}

\begin{itemize}
  \item רישום: דמי רישום חד-פעמיים של 25 דולר; נדרש חשבון Google תקין;
  \item חומרים: אייקון אפליקציה (512×512), צילומי מסך (לפחות 2), תיאור התכונות וקישור למדיניות פרטיות (חובה כשאוספים נתונים);
  \item טופס אבטחת נתונים: מצהירים ביושר על הנתונים שנאספים – ספר חשבונות ללא חיבור לרשת יכול להצהיר "אפס איסוף";
  \item ביקורת: חשבונות חדשים בהעלאה הראשונה עוברים לרוב בדיקה אנושית, מכמה שעות ועד כמה ימים; אם דוחים – מתקנים לפי ההערות.
\end{itemize}

לתשומת לב משתמשים בסין: Google Play דורש סביבת רשת מיוחדת, ובסין היבשתית ההפצה נעשית בדרך כלל דרך החנויות המקומיות (Huawei, Xiaomi, MyApp וכדומה). תהליך העלאה לחנות מקומית דומה, אך לכל חנות טעם ביקורת משלה – אותו AAB/APK אפשר להפיץ בכמה חנויות. לצורך תרגול – לעבור בהצלחה חנות אחת כבר ניצחון.

\booksection{ה' איך מכינים את חומרי החנות: הרושם הראשון שווה הורדות}

עמוד החנות הוא הוויטרינה שהמשתמש רואה ב“טיול בחנות”. שלושה חומרים קובעים את הרושם הראשוני:

\begin{itemize}
  \item אייקון: ‏512×512, אל תסתפקו בצילום מסך. אין כישורי עיצוב? צרו עם כלי AI צורה גיאומטרית פשוטה; האייקון עובר רק אם מזהים אותו במבט אחד גם כשהוא מוקטן על מסך הבית של הטלפון;
  \item צילומי מסך: לפחות 2, עדיף 4-6, שמציגים את תרחישי השימוש המרכזיים; לכל אחד שורת טקסט בשפה פשוטה ("רושמים עסקה בשנייה" טובה פי עשרה מ"פתרון יעיל לניהול חשבונות");
  \item תיאור: המשפט הראשון מסביר "מה זה + למי זה מיועד"; אל תערמו "הכי מקצועי, הכי חזק" – משקל החיפוש בחנות נקבע בעיקר לפי מילות מפתח ושיעור המרה להורדות, וכנות יעילה יותר מהגזמה.
\end{itemize}

\begin{tcbinfo}{טיפ: עשו A/B והסתכלו על המתחרים}
לפני העלייה, סיירו בעשירייה הפותחת של אפליקציות דומות בחנות והשוו צילומי מסך: איזה סגנון אייקון יש להן? מה נקודת המכירה בצילומים? איך נוסח הכותרת? חיקוי אינו העתקה – לומדים את המתכונת ומחליפים לתוכן משלכם; זהו השיעור המהיר ביותר להשקה קרה של כל מוצר.

\end{tcbinfo}

\booksection{ה' רשימה קצרה לאופטימיזציית ביצועים ודיווח קריסות}

לפני שהאפליקציה עולה לחנות, שווה להקדיש לכל אחד משני הדברים האלה חצי שעה:

\booksubsection{ארבעה סעיפים לבדיקת ביצועים עצמית}

\begin{enumerate}
  \item רשימות ב-LazyColumn: לעולם אל תסדרו רשימה ארוכה ב-Column רגיל – טעינה עצלה היא השיעור הראשון בביצועי אנדרואיד;
  \item אל תדחפו תמונות מקור בגודל מלא: השתמשו בספריית התמונות Coil לשינוי קנה מידה ומטמון אוטומטיים; תמונת מקור אחת של 5MB מספיקה כדי לחנוק מכשיר חלש;
  \item הימנעו מעבודה כבדה בת'רד הראשי: מסד נתונים, רשת וקבצים גדולים – הכול בתוך קורוטינות;
  \item מדדו עם Profiler: כלי ניתוח הביצועים המובנה ב-Android Studio מציג זיכרון, CPU וצריכת סוללה – קודם מודדים, אחר כך מייעלים; אל תסתמכו על ההרגשה.
\end{enumerate}

\booksubsection{דיווח קריסות}

זה שלא קרס אצלכם בבדיקות לא אומר שלא יקרוס בטלפון של משתמשים. לפני העלייה חברו את Firebase Crashlytics (הרשמי של גוגל, חינמי): כמה שורות הגדרה, ובכל קריסת משתמש נשלחת אוטומטית המחסנית – בממשק הניהול תראו “איזה דגם, איזו גרסה, ובאיזו שורה קרס”. עלייה לחנות בלי ניטור קריסות שווה נהיגה בעיניים עצומות.

\begin{tcbpractice}{הגיע הזמן להתנסות: "למסור" את ספר החשבונות פעם אחת}
(1) הוסיפו 3 בדיקות יחידה ללוגיקה כמו parseAmount; (2) תפסו במו ידיכם ב-Logcat קריסה שיצרתם בעצמכם ותקנו אותה; (3) צרו APK עם חתימת release ושלחו לחבר להתקנה (זהו בדיוק טקס “היצירה”); (4) אם נשאר כוח – רשמו חשבון מפתחים והגישו באמת את ספר החשבונות לחנות. את שלושת הראשונים אפשר להשלים עוד הערב.

\end{tcbpractice}

\begin{tcbsummary}{סיכום הפרק}
\begin{itemize}
  \item שלושת כלי הניפוי: Logcat ללוגים, נקודות עצירה למשתנים, Layout Inspector למסך
  \item קריאת קריסה: מלמטה למעלה במחסנית, מוצאים את השורה הראשונה עם שם החבילה שלכם
  \item שלוש שכבות בדיקות: יחידה (שניות – מתחילים בזה)/מכשיר/ידנית
  \item הפצה מחייבת חתימה: גבו את ה-keystore ואת הסיסמה; אם אבדו – לאפליקציה "מחליפים זהות"
  \item ששת שלבי העלייה לחנות: רישום חומרים העלאת AAB טופס אבטחת נתונים ביקורת פרסום
\end{itemize}

\end{tcbsummary}
